Utilizando a combinação de dados de LiDAR, câmeras, radar e GNSS é possível criar um modelo situacional robusto, mesmo em ambientes de risco. A principal aplicação é evitar colisões com detritos espaciais ou realizar manobras orbitais em ambientes congestionados com uma quantidade elevada de detritos ou outros equipamentos espaciais.

\subsection{Fusão de sensores}

A fusão de sensores consiste na combinação de dados sensoriais ou de informações derivadas desses dados para gerar uma representação mais precisa do ambiente do processo, dessa forma a abordagem é melhorada e com isso, eleva a robustez dos sistemas, ampliando a cobertura espacial e temporal, portanto reduz ambiguidades e incertezas e aprimora a resolução das medições \cite{elmenreichDissertation}.

A abordagem baseada em eventos temporizados (time-triggered) é uma estrutura adequada para a implementação de fusão de sensores em sistemas de tempo real, devido ao seu comportamento determinístico. A integração dessa técnica em um framework de tempo-triggered permite a construção de sistemas confiáveis e eficientes em termos de recursos. Para que essa aplicação seja viável, o sistema pode ser dividido em três níveis, sendo que o primeiro nível (transdutor), compreende os sensores e atuadores, já o segundo nível, denominado fusão e disseminação, coleta as medições, executa a fusão dos dados e distribui as informações processadas aos atuadores, por fim, no terceiro nível (controle), um programa de controle processa os dados do nível de fusão e toma decisões com base nas informações ambientais recebidas. Essa estrutura permite que aplicações complexas sejam fragmentadas em subsistemas menores e gerenciáveis.

Ainda conforme explica \cite{elmenreichDissertation}, no nível do transdutor, um algoritmo de filtragem, como o filtro de Kalman, podem ser utilizado para melhorar a qualidade dos dados de um único sensor por meio de medições sucessivas. Outra abordagem, aplicada no nível de fusão e disseminação, combina os dados provenientes de múltiplos sensores para obter informações mais consistentes e precisas. Entre os métodos propostos, destacam-se o algoritmo de ponderação sistemática de confiança e o algoritmo de grade de certeza robusta, desenvolvido para aplicações específicas.

\section{Filtro de Kalman}

Conforme \cite{kalman1960}, o filtro de Kalman é um método utilizado para estimar estados de um sistema dinâmico a partir de medições ruidosas, ele combina previsões do modelo do sistema com medições observadas para obter uma estimativa otimizada, levando em consideração as incertezas associadas a cada informação. Este filtro foi desenvolvido por Kalman e Bucy em 1960, e desde então tem sido aplicado em diversas áreas, como engenharia aeroespacial, robótica e navegação.

A abordagem do filtro baseia-se em dois estágios principais: previsão e atualização. No estágio de previsão, a estimativa do estado do sistema é calculada a partir do modelo matemático, levando em conta a relação entre os estados anteriores e as entradas aplicadas, em seguida, ocorre a atualização, na qual os dados observados são incorporados à estimativa anterior, ajustando-a de acordo com um fator de ponderação chamado ganho de Kalman. Esse ganho determina o peso dado às medições recentes em relação à previsão do modelo, garantindo que medições mais confiáveis tenham maior influência no resultado.

Um dos principais benefícios do filtro de Kalman é sua capacidade de lidar com ruído e incertezas nos dados, ele assume que tanto o sistema quanto as medições estão sujeitos a erros estatísticos e modela essas incertezas utilizando distribuições de probabilidade, portanto, ele consegue fornecer estimativas mais confiáveis ao longo do tempo, mesmo quando há falhas momentâneas ou medições inconsistentes. Além disso, a sua implementação é não requer poder computacional elevado, conforme explica \cite{gan2001}, permitindo sua aplicação em sistemas em tempo real, como navegação de veículos autônomos, rastreamento de objetos e até sistemas de navegação como utilizado pelo Boeing 777 \cite{cipra1993}.

Ao longo dos anos, o modelo original do filtro de Kalman foi expandido para atender a aplicações mais complexas, um exemplo disso é o filtro para sistemas não lineares, a versão conhecida como Filtro de Kalman Estendido (EKF) introduz uma linearização das equações do sistema para que possam ser processadas dentro da estrutura original do filtro \cite{welch2001}.

O filtro de Kalman tem sido amplamente empregado na fusão de sensores, como explica \cite{li2015}, ao combinar informações de múltiplas fontes para produzir estimativas mais precisas. Em sistemas de navegação, por exemplo, ele é utilizado para integrar dados de sistemas de posicionamento global (GNSS) e sensores inerciais (INS), compensando limitações de cada tecnologia individualmente. Quando sinais GNSS são temporariamente perdidos devido a bloqueios por edifícios ou interferências, o filtro de Kalman pode manter a estimativa da posição com base nas medições inerciais, reduzindo erros acumulativos.

